\section{Related work}\label{sec:related}

\paragraph{RFM segmentation} RFM summarises a customer's history by the time since the last purchase, the number of purchases and the amount spent \citep{hughes1994}. Its predictive value is well established: \citet{fader2005} show that RFM is a sufficient summary for customer-lifetime-value projection under standard purchase models, using the CDNOW data that we also use. Segmentation practice scores each dimension by quintile and either reads segments off score combinations or clusters customers, most often with K-means; \citet{chen2012} introduced the Online Retail data with an RFM/K-means case study, and K-means, hierarchical clustering and Gaussian mixtures remain the usual choices \citep{wedel2000}. Cluster-based segments, however, are described only by centroids and must be labelled after the fact, which is the gap FCA addresses.

\paragraph{Fuzzy RFM} Fuzziness has entered RFM work in two ways. \citet{carrasco2019} replace crisp RFM scores by 2-tuple fuzzy linguistic values before K-means segmentation, addressing the imprecision of ordinal scores. \citet{marindiaz2025} cluster RFM features with fuzzy C-means \citep{bezdek1981}, map them to a 2-tuple linguistic scale and add explainable-AI post-processing. In both, the segments are still clusters; fuzziness softens the inputs or the assignment, but segments do not carry explicit, exact descriptions or an order between them. Our work keeps FCA's explicit, ordered segments and makes the context itself fuzzy.

\paragraph{Formal concept analysis} FCA derives a lattice of concepts (extent, intent) from a binary object--attribute context \citep{ganter1999}. Surveys document its use for knowledge discovery, information retrieval and ontology engineering \citep{poelmans2013app,poelmans2013mod,singh2016}. Two practical issues recur: the number of concepts grows quickly, which motivates pruning by support, stability \citep{kuznetsov2007} or other indices, and binary contexts require discretising numeric data. Association rules between concepts are usually read along lattice covers \citep{luxenburger1991}. \citet{rungruang2024} applied FCA to quintile RFM scores and showed that the lattice gives an interpretable, hierarchical segmentation of Online Retail II customers, compared with K-means and hierarchical clustering by cluster-validity indices.

\paragraph{Fuzzy formal concept analysis} Fuzzy FCA generalises the binary incidence relation to degrees in a residuated lattice \citep{burusco1994,pollandt1997,belohlavek2002,belohlavek2004}; surveys review its variants and applications \citep{poelmans2014,alwersh2022}. A recurring issue, noted by \citet{alwersh2022}, is the much larger number of fuzzy concepts, which makes concept reduction necessary for applications. We address it with a scale-aware support notion (core support) and a summary with an approximation guarantee, and we show why the common sigma-count support is inadequate under the Łukasiewicz residuum.

\paragraph{Positioning} To our knowledge, fuzzy formal concepts have not been used for RFM segmentation, and RFM-FCA has not been evaluated out of sample, against clustering baselines on prediction, targeting and stability, or across several datasets. This paper does all three.
